% response-letter: tcolorbox styles for the response divs.
%
% This file owns the GEOMETRY and TYPOGRAPHY of the boxes only. The color
% palette (rlComment, rlReply, rlChanges, rlQuoted) is injected into the
% preamble by response-letter.lua from the `response-letter: colors:`
% metadata — change colors in the YAML, not here. tcolorbox resolves colors
% at use time, so the definition order does not matter.
\usepackage[svgnames]{xcolor}
\usepackage[most]{tcolorbox}
\usepackage{caption}
\usepackage{pifont}

\newtcolorbox{CommentBox}{
  breakable, enhanced,
  colback=rlComment!5, colframe=rlComment!60,
  boxrule=0pt, leftrule=2.5pt, arc=1pt,
  left=7pt, right=7pt, top=6pt, bottom=6pt,
  before skip=10pt, after skip=10pt,
  fontupper=\itshape
}

\newtcolorbox{ReplyBox}{
  breakable, enhanced,
  colback=rlReply!4, colframe=rlReply!45,
  boxrule=0pt, leftrule=2.5pt, arc=1pt,
  left=7pt, right=7pt, top=6pt, bottom=6pt,
  before skip=10pt, after skip=10pt
}

\newtcolorbox{ChangesBox}{
  breakable, enhanced,
  colback=rlChanges!12, colframe=rlChanges!80,
  boxrule=0pt, leftrule=2.5pt, arc=1pt,
  left=7pt, right=7pt, top=6pt, bottom=6pt,
  before skip=10pt, after skip=10pt
}

\newtcolorbox{QuoteBox}{
  breakable, enhanced,
  colback=rlQuoted!8, colframe=rlQuoted,
  boxrule=0pt, leftrule=3pt, arc=0pt,
  left=9pt, right=7pt, top=6pt, bottom=6pt,
  before skip=8pt, after skip=8pt
}
